\section{Literature Review}
\noindent
Digital signal processing (DSP) refers to the manipulation of signals represented
in discrete form, such as those stored or processed by computers. In computer music, DSP 
provides the theoretical and practical foundation for digitally modelling and recreating the 
behaviour of analog audio systems, including oscillators and filters. However, accurately 
translating analog oscillators and filters into the digital domain presents several technical 
challenges. Specifically, the choice of oscillator anti-aliasing techniques and digital filter 
architectures introduces different considerations that influence the accuracy, stability, and 
computational efficiency of the resulting audio. These considerations become particularly 
important in real-time browser-based synthesis, where audio fidelity must be balanced against 
computational constraints.

This section examines these challenges through a review of existing 
literature, focusing on the methods or approaches used or proposed to address them.

\subsection{Digital Oscillators}
\noindent
Analog oscillators generate the fundamental voltage signals used in analog synthesisers,
while digital oscillators aim to reproduce equivalent waveforms in the discrete-time domain.
Common waveforms include the sine, sawtooth, square and triangle waves, each with distinct
harmonic characteristics that contribute to its perceived timbre.

A naive digital oscillator generates waveforms directly from mathematical functions or
phase-based calculations. A sawtooth, for example, can be produced with a phase accumulator,
which advances its phase on every sample by an amount determined by the oscillator frequency
and the sample rate. The phase is mapped to an amplitude range and, once it reaches its
maximum, wraps back to the minimum, producing the characteristic sawtooth shape. This abrupt
jump is known as a \textit{discontinuity}, and it introduces infinitely many high-frequency 
harmonics. The
Nyquist--Shannon sampling theorem states that a signal can only be represented accurately if
it contains no components above the Nyquist frequency, which is half the sampling frequency.
A sawtooth always contains harmonics above this limit, so they cannot be represented and
instead fold back into the audible range as inharmonic distortion known as aliasing.

The problem worsens as the fundamental frequency rises. This applies to any waveform with
discontinuities, such as the sawtooth and also the square wave, because a greater proportion of
their harmonics then lies above the Nyquist frequency.

\subsection{Anti-Aliasing Techniques}
\noindent
Anti-aliasing refers to techniques that reduce the effect of aliasing. One approach is the
Band-Limited Step (BLEP) method. The idea is to replace the abrupt discontinuity in the naive
waveform with a band-limited approximation, which reduces the high-frequency content introduced
by the discontinuity and so improves the accuracy of the eventual reconstruction.

An ideal step is the integral of an impulse, and an ideal band-limited step is therefore
the integral of a band-limited impulse, typically a windowed sinc function\cite{BLEP}. 
Rather than synthesising the entire waveform from band-limited components, the BLEP method 
computes only the \textit{residual}---the difference between the band-limited step and the naive 
(non-band-limited) step---and applies it as a correction to the naive waveform in the immediate 
vicinity of each discontinuity, reducing the effect of aliasing.

A widely used variant is the polynomial BLEP (PolyBLEP), which approximates the residual
with a low-order polynomial instead of a table of sinc-derived values\cite{PolyBLEP}. For a 
normalised phase
$t \in [0, 1)$ and a per-sample phase increment $\Delta t = f_0 / f_s$, where $f_0$ is the
oscillator frequency and $f_s$ the sampling frequency, a two-sample PolyBLEP residual is given by
\begin{equation}
  r(t) =
  \begin{cases}
    2\dfrac{t}{\Delta t} - \left(\dfrac{t}{\Delta t}\right)^{2} - 1,
      & 0 \le t < \Delta t, \\[2ex]
    \left(\dfrac{t - 1}{\Delta t}\right)^{2} + 2\dfrac{t - 1}{\Delta t} + 1,
      & 1 - \Delta t < t < 1, \\[2ex]
    0, & \text{otherwise.}
  \end{cases}
\end{equation}
The corrected sawtooth is then obtained by subtracting this residual from the ideal but naive 
sawtooth waveform $s(t)$,
also smoothing the wrap-around over the samples adjacent to the discontinuity. The naive sawtooth 
wave is chosen to be $s(t)=2t-1$ for simplicity of residual calculation and thus implementation 
of PolyBLEP, because it reaches its minimal extreme of $s(t)=-1$ at $t=0$, and the maximum extreme 
of $s(t)=1$ at $t=1$; the discontinuities exist at integer intervals.

A square wave
may be treated as the sum of two sawtooth-like discontinuities, one rising and one falling, so
the same residual is applied at both transitions (with opposite sign). A triangle wave, being continuous, has no step discontinuities,
but its first derivative is discontinuous; the slope discontinuity can be corrected with an
analogous polynomial band-limited ramp (PolyBLAMP) residual\cite{PolyBLAMP}

The principal advantages of PolyBLEP and PolyBLAMP are their low computational cost, the lack of 
lookup tables,
and suitability for real-time use with modulated frequency, since the correction depends
only on the current phase and phase increment. Its limitation is that a low-order polynomial
only approximates the ideal band-limited step, so some residual aliasing remains, particularly
at high fundamental frequencies, where it is attenuated rather than eliminated.

Alternative approaches offer different trade-offs. Oversampling reduces aliasing by generating 
waveforms at higher sampling rates before filtering and decimation, but increases computational cost
\cite{oversampling}. Wavetable synthesis uses precomputed band-limited waveforms to suppress harmonics 
above the Nyquist frequency, but requires additional memory and table selection, reducing flexibility 
especially during dynamic modulation\cite{wavetable}.
\begin{comment}
Alternative
approaches offer other trade-offs
between aliasing suppression, computational cost and flexibility. Oversampling improves anti-aliasing
by generating the waveform at a higher sampling rate, thereby pushing aliasing components further above 
the audible frequency range before low-pass filtering and decimation. 
However, its limitation is the increased computational cost associated with generating and 
filtering additional samples from the increased sampling rate\cite{oversampling}. Wavetable synthesis 
improves aliasing suppression by using precomputed band-limited waveforms with progressively 
reduced harmonic content at higher frequencies, ensuring that only harmonics below the Nyquist 
frequency are used. However, its limitation is the memory required to store 
multiple tables and 
the need for table selection and interpolation, which can introduce additional complexity and 
may reduce flexibility when the waveform is dynamically modulated\cite{wavetable}.
\end{comment}


